\input{../tex-inputs/latex-korrekturansicht-vorspann}

\section[Arthur Schnitzler an Gustav Schwarzkopf, 2. 1. 1929]{L04352 Arthur Schnitzler an Gustav Schwarzkopf, 2. 1. 1929}
\nopagebreak\mylabel{L04352v}
\rehead{ }\normalsize\beginnumbering\briefempfaengerindex{Schwarzkopf, Gustav@\textsc{Schwarzkopf, Gustav}!zzzSchnitzler, Arthur@\emph{von Arthur Schnitzler}!1929-01-021@{2. 1. 1929}|(be}
\toendnotes[C]{\smallbreak\pagebreak[2]}
\correspDesc{Versand  durch Arthur Schnitzler am 2. 1. 1929 in Berlin
\newline{}Übermittlung  am 2. 1. 1929 in Berlin
\newline{}Erhalt  durch Gustav Schwarzkopf im Zeitraum [3. 1. 1929 – 7. 1. 1929?] in Wien}\toendnotes[C]{\smallbreak}
\Standort{DLA, A:Schnitzler, HS.1985.1.1897.}
\physDesc{Bildpostkarte, 110 Zeichen
\newline{}Handschrift: Bleistift, lateinische Kurrent
\newline{}Versand: 1) Stempel: »\nobreak{}\oindex{Berlin@\textbf{Berlin}, \emph{Hauptstadt}|pwk}{[}Be{]}rlin SW 11, 2 1 2\textcolor{gray}{9}, 10–11N\nobreak{}«.   2) Stempel: »\nobreak{}\textcolor{brown}{Grüne Woche Berlin}\orgindex{Grüne Woche@Grüne Woche|pw} 26. 1. – 3. 2. 1929\nobreak{}«. }\pstart{}{\pb}Hrn Gustav Schwarzkopf\pend{}\pstart{}\textcolor{pink}{Wien I}\oindex{I., Innere Stadt@\textbf{I., Innere Stadt}, \emph{Verwaltungsgebiet}|pw}{}\ledrightnote{\textcolor{pink}{I., Innere Stadt}}\pend{}\pstart{}\textcolor{pink}{Tiefer Graben 17}\oindex{Wien@\textbf{Wien}!I., Innere Stadt@\textbf{I., Innere Stadt}!Tiefer Graben 17@\textbf{Tiefer Graben 17}, \emph{Wohngebäude}|pw}{}\ledrightnote{\textcolor{pink}{Tiefer Graben 17}}\pend{}{\bigskip}
\pstart
           {\pb}\textcolor{pink}{\textcolor{gray}{\textbf{Hotel Esplanade, Berlin}}}\oindex{Hotel Esplanade [Berlin]@\textbf{Hotel Esplanade [Berlin]}, \emph{Hotel}|pw}{}\ledrightnote{\textcolor{pink}{Hotel Esplanade [Berlin]}}\hfill \textcolor{gray}{\textbf{Vorhalle des Winter-Restaurants}}\pend
           \vspace{1em}
\pstart
           \raggedleft{}{\pb}\textcolor{pink}{Berlin}\oindex{Berlin@\textbf{Berlin}, \emph{Hauptstadt}|pw}{}\ledrightnote{\textcolor{pink}{Berlin}}{ }2. 1. 929\pend
           \vspace{0.5em}
\pstart
           Herzliche Grüße und auf baldgs Wiedersehn\pend
           
\pstart
           Ihr{\\[\baselineskip]}\spacefill\mbox{ArSchn}\pend
           \leftskip=0em{}\selectlanguage{ngerman}\endnumbering\briefempfaengerindex{Schwarzkopf, Gustav@\textsc{Schwarzkopf, Gustav}!zzzSchnitzler, Arthur@\emph{von Arthur Schnitzler}!1929-01-021@{2. 1. 1929}|)be}\mylabel{L04352h}
\begin{anhang}
\end{anhang}\newcommand{\dateiname}{L04352}\newcommand{\titel}{Arthur Schnitzler an Gustav Schwarzkopf, 2. 1. 1929}\newcommand{\editorInnen}{Herausgegeben von Jahnke, SelmaMüller, Martin Anton}\input{../tex-inputs/latex-korrekturansicht-abspann}
